\section{Linear scalar perturbations and CLASS variable map}\label{app:linear}
The quadratic derivation is organized so that the physical CLASS scalar combination is exposed before implementation. In the notation used during the variation,
\begin{equation}
\theta_s=\frac{\varphi}{\Q},\qquad
\alpha=\frac{\chi}{\Q}-\theta_s,
\end{equation}
so
\begin{equation}
\varphi+\Q\alpha=\chi.
\end{equation}
The quadratic contribution from the nonseparable operator therefore reduces to
\begin{equation}
\Delta\mathcal L^{(2)}=-\Gamma(\Q)\frac{k^2}{a^2}\chi^2
\end{equation}
in the physical scalar basis.

The generic direct source entering the $E$ equation contains
\begin{equation}
K_Q\chi-\Q F_{,\Y}\chi-(2-K_B)B.
\end{equation}
The current identity gives
\begin{equation}
K_Q=\Q F_{,\Y}=\frac{I_0}{a^3},
\end{equation}
so the direct $\chi$ source cancels exactly and the corresponding reduced equation contains
\begin{equation}
K_B(\dot E+HE)=-(2-K_B)B
\end{equation}
for this direct source contribution. The numerical source therefore inserts an analytic zero rather than a floating-point subtraction.

The mixed derivative $F_{,\Y\Q}=\Gamma_{,\Q}$ first contributes at cubic order. The kinematic definitions of the AeST density, velocity, $\chi$, anisotropic stress, and composite perturbation variables are unchanged by the nonseparable operator. In the auxiliary acceleration sector, the nonzero-$k$ scalar perturbations of the auxiliary field and gauge-fixing variables are eliminated algebraically, so the physical propagating scalar basis remains that of the scalar--vector gravitational sector after the metric constraints are solved.
